% CP-VI-0048
% target_chapter: VI/48 — Exact Sequences and Cohomology
% source_unit_id: PVI-RECOVER-A99C7658F3-L000557
% source: Downloads/theory-of-algebraic-geometry-2.tex L557-L836
% solution_provenance: SOURCE_EXPLICIT_SOLUTION

\subsection*{Problem 9. The Functor \texorpdfstring{\(\Gamma(U,-)\)}{Gamma(U,-)} is Left Exact}

Let \(U\subseteq X\) be an open subset. Recall that for a sheaf \(\mathcal{F}\),
\[
\Gamma(U,\mathcal{F})=\mathcal{F}(U)
\]
is the group of sections of \(\mathcal{F}\) over \(U\).

We want to show that the functor
\[
\Gamma(U,-)
\]
is left exact.

That is, if
\[
0\longrightarrow \mathcal{F}^{\prime}
\xrightarrow{i}
\mathcal{F}
\xrightarrow{p}
\mathcal{F}^{\prime\prime}
\]
is an exact sequence of sheaves, then the induced sequence of abelian groups
\[
0\longrightarrow \Gamma(U,\mathcal{F}^{\prime})
\xrightarrow{\Gamma(U,i)}
\Gamma(U,\mathcal{F})
\xrightarrow{\Gamma(U,p)}
\Gamma(U,\mathcal{F}^{\prime\prime})
\]
is exact.

Equivalently, we must show two things:

\[
\Gamma(U,i):\Gamma(U,\mathcal{F}^{\prime})\longrightarrow \Gamma(U,\mathcal{F})
\]
is injective, and
\[
\operatorname{im} \Gamma(U,i)=\ker \Gamma(U,p).
\]

\subsection{Step 1: Injectivity}

Let
\[
s\in \Gamma(U,\mathcal{F}^{\prime})
\]
and suppose
\[
\Gamma(U,i)(s)=0
\]
in \(\Gamma(U,\mathcal{F})\). In simpler notation, this means
\[
i(s)=0.
\]

Passing to germs, for every \(x\in U\), we get
\[
i_x(s_x)=0.
\]
Because the original sequence of sheaves is exact, the stalk map
\[
i_x:\mathcal{F}^{\prime}_x\longrightarrow \mathcal{F}_x
\]
is injective. Therefore
\[
s_x=0
\]
for every \(x\in U\).

A section of a sheaf is zero if and only if all of its germs are zero. Hence
\[
s=0.
\]
Therefore
\[
\Gamma(U,i)
\]
is injective.

Thus the sequence
\[
0\longrightarrow \Gamma(U,\mathcal{F}^{\prime})
\longrightarrow \Gamma(U,\mathcal{F})
\]
is exact.

\subsection{Step 2: The image is contained in the kernel}

Let
\[
s\in \Gamma(U,\mathcal{F}^{\prime}).
\]
Then
\[
\Gamma(U,p)(\Gamma(U,i)(s))
=
p(i(s)).
\]
Since the sheaf sequence is exact, we have
\[
p\circ i=0.
\]
Therefore
\[
p(i(s))=0.
\]
Hence
\[
\Gamma(U,i)(s)\in \ker \Gamma(U,p).
\]
So
\[
\operatorname{im} \Gamma(U,i)\subseteq \ker \Gamma(U,p).
\]

\subsection{Step 3: The kernel is contained in the image}

Now let
\[
t\in \Gamma(U,\mathcal{F})
\]
and suppose that
\[
\Gamma(U,p)(t)=0.
\]
Equivalently,
\[
p(t)=0
\]
in \(\Gamma(U,\mathcal{F}^{\prime\prime})\).

Passing to germs, for every \(x\in U\), we have
\[
p_x(t_x)=0.
\]
Hence
\[
t_x\in \ker(p_x).
\]

Because the original sequence of sheaves is exact, for every \(x\in U\),
\[
\ker(p_x)=\operatorname{im}(i_x).
\]
Therefore for every \(x\in U\), there exists a germ
\[
s_x\in \mathcal{F}^{\prime}_x
\]
such that
\[
i_x(s_x)=t_x.
\]

By the definition of a germ, for every \(x\in U\), there exists an open neighbourhood
\[
V_x\subseteq U
\]
and a section
\[
s_{V_x}\in \mathcal{F}^{\prime}(V_x)
\]
such that
\[
i(s_{V_x})=t|_{V_x}.
\]

Now take two such open sets \(V_x\) and \(V_y\). On their overlap \(V_x\cap V_y\), we have
\[
i(s_{V_x}|_{V_x\cap V_y})
=
t|_{V_x\cap V_y}
=
i(s_{V_y}|_{V_x\cap V_y}).
\]
Since \(i\) is injective, it follows that
\[
s_{V_x}|_{V_x\cap V_y}
=
s_{V_y}|_{V_x\cap V_y}.
\]

Thus the sections \(s_{V_x}\) agree on overlaps. Since \(\mathcal{F}^{\prime}\) is a sheaf, the gluing axiom gives a unique section
\[
s\in \mathcal{F}^{\prime}(U)=\Gamma(U,\mathcal{F}^{\prime})
\]
such that
\[
s|_{V_x}=s_{V_x}
\]
for all \(x\in U\).

Then on each \(V_x\),
\[
i(s)|_{V_x}
=
i(s|_{V_x})
=
i(s_{V_x})
=
t|_{V_x}.
\]
Hence, by the locality axiom for the sheaf \(\mathcal{F}\),
\[
i(s)=t.
\]
Therefore
\[
t\in \operatorname{im} \Gamma(U,i).
\]
So
\[
\ker \Gamma(U,p)\subseteq \operatorname{im} \Gamma(U,i).
\]

Combining the two inclusions, we obtain
\[
\ker \Gamma(U,p)=\operatorname{im} \Gamma(U,i).
\]

Therefore the sequence
\[
0\longrightarrow \Gamma(U,\mathcal{F}^{\prime})
\longrightarrow \Gamma(U,\mathcal{F})
\longrightarrow \Gamma(U,\mathcal{F}^{\prime\prime})
\]
is exact.

Hence
\[
\boxed{
	\Gamma(U,-)\text{ is left exact.}
}
\]

\subsection[Important warning: Gamma(U,-) is not necessarily exact]{Important warning: \texorpdfstring{\(\Gamma(U,-)\)}{Gamma(U,-)} is not necessarily exact}

The functor \(\Gamma(U,-)\) is left exact, but it is not necessarily exact.

Indeed, suppose
\[
0\longrightarrow \mathcal{F}^{\prime}
\longrightarrow \mathcal{F}
\longrightarrow \mathcal{F}^{\prime\prime}
\longrightarrow 0
\]
is a short exact sequence of sheaves. Applying \(\Gamma(U,-)\), we always get an exact sequence
\[
0\longrightarrow \Gamma(U,\mathcal{F}^{\prime})
\longrightarrow \Gamma(U,\mathcal{F})
\longrightarrow \Gamma(U,\mathcal{F}^{\prime\prime}).
\]
However, the map
\[
\Gamma(U,\mathcal{F})\longrightarrow \Gamma(U,\mathcal{F}^{\prime\prime})
\]
need not be surjective.

The reason is that surjectivity of a morphism of sheaves means local liftability, not necessarily global liftability.

In other words, if
\[
t\in \Gamma(U,\mathcal{F}^{\prime\prime}),
\]
then locally on \(U\), one can find sections of \(\mathcal{F}\) mapping to \(t\). But these local lifts may fail to glue to a global section of \(\mathcal{F}(U)\).

Thus, in general,
\[
\Gamma(U,-)
\]
preserves injectivity and kernels, but not necessarily surjectivity.

\[
\boxed{
	\Gamma(U,-)\text{ is left exact, but not necessarily exact.}
}
\]
